\section{Related work}
\label{sec:related}

\subsection{Spatial cognition framework}

The psychology of spatial knowledge acquisition provides a foundational taxonomy for understanding geographic reasoning \citep{siegel1975development}. This landmark-route-survey framework distinguishes three hierarchical but complementary forms of spatial knowledge:

\begin{itemize}[leftmargin=*, itemsep=1pt]
    \item \textbf{Landmark knowledge:} Recognition and recall of salient environmental features (buildings, intersections, natural features) without explicit encoding of metric distances or spatial relations between them.
    \item \textbf{Route knowledge:} Ordered sequences of actions or paths connecting landmarks, represented as procedural, egocentric knowledge tied to specific traversal experiences (e.g., ``turn left at the church, then walk straight to the bridge'').
    \item \textbf{Survey knowledge:} Allocentric, map-like understanding of space encoding global layout, metric distances, and directional relationships, enabling flexible reasoning such as shortcut discovery, detour planning, and novel route inference.
\end{itemize}

These knowledge types are acquired through different learning modalities and exhibit distinct strengths and limitations. \citet{thorndyke1982differences} demonstrate that navigation-based learning preferentially supports route knowledge, whereas map-based learning facilitates survey knowledge. Subsequent work has shown that cognitive maps---internal representations of spatial environments---develop progressively from landmark to route to survey knowledge \citep{montello1998new}.

\paragraph{Extending the framework for LLM evaluation.}
We adopt and extend this framework for \dsName by adding a fourth category: \textbf{Geometric (G)} knowledge, representing pure mathematical operations on coordinates that require no geographic world knowledge. This extension serves as a control condition, isolating computational ability from world-knowledge integration. For example, computing the Haversine distance between two coordinate pairs is purely geometric, while identifying which city lies at those coordinates requires landmark knowledge. This decomposition enables \dsName to diagnose whether model failures stem from missing geographic facts or insufficient mathematical reasoning---a distinction obscured in existing benchmarks that conflate these capabilities.

\subsection{Geospatial benchmarks}

Evaluating large language models on geospatial tasks has emerged as an active research area, with benchmarks spanning geographic knowledge assessment, spatial reasoning, coordinate-based computation, and tool-augmented workflows.

\paragraph{Geographic knowledge and bias.}
Several benchmarks evaluate LLM factual geographic knowledge and expose systematic biases. WorldBench \citep{moayeri2024worldbench} tests country-level indicator recall across 20 LLMs, revealing 1.5$\times$ higher error rates for Sub-Saharan Africa versus North America and nearly tripled error rates between the lowest and highest income countries. \citet{manvi2024geographically} demonstrate systematic biases against low-socioeconomic regions across multiple model families. GeoLLM \citep{manvi2024geollm} shows LLMs embed substantial geospatial knowledge (population density, income levels) extractable through fine-tuning with auxiliary OpenStreetMap data, achieving 70\% improvement over naive coordinate prompting. GeoGLUE \citep{li2023geoglue} evaluates geographic language understanding across textual similarity and entity alignment tasks. LocalBench \citep{li2024localbench} extends geographic evaluation to county-level knowledge in the United States, testing fine-grained local facts beyond country-level statistics.

\paragraph{Spatial reasoning.}
Spatial reasoning has received increasing attention alongside broader interest in general-purpose reasoning for LLMs. Early benchmarks focus on abstract, grid-based, or shape-centric reasoning. \citet{yamada2024evaluating} show that LLMs possess basic understanding of common shapes like squares and triangles but struggle with complex shapes like rings and hexagons. SpatialEval \citep{wang2024spatial} evaluates multimodal spatial understanding across navigation, grid reasoning, and map interpretation tasks, revealing substantial performance degradation as spatial complexity increases. For vision-language models, SpatialBench \citep{xu2025spatialbench} focuses on 3D and metric spatial reasoning, exposing persistent failures in depth, proximity, and relative positioning, while 3DSRBench \citep{ma20243dsrbench} provides 2,772 manually annotated question-answer pairs across 12 question types covering height, location, orientation, and multi-object reasoning.

More aligned with geographical reasoning, GeoGramBench \citep{luo2025geogrambench} tests translation of geometry code to spatial representations, with frontier models achieving $<$50\% accuracy at higher abstraction levels. MapEval \citep{dihan2025mapeval} assesses map-based reasoning through 700 questions spanning 180 cities across textual, API-based, and visual modalities, where no model surpasses 67\% accuracy and all lag over 20\% behind human performance. GeoEval \citep{zhang2024geoeval} evaluates geometry problem-solving with both diagrams and text, revealing that even GPT-4V achieves only 55.6\% on plane geometry problems. For geoparsing, benchmarks like LGL \citep{lieberman2010geotagging} and WikToR \citep{gritta2018wikor} evaluate text-to-coordinate mapping.

\paragraph{Tool-use geospatial tasks.}
A separate line of work evaluates LLMs' ability to orchestrate external geographic information system (GIS) tools. GeoBenchX \citep{krechetova2025geobenchx} benchmarks over 200 multi-step tasks involving 23 GIS tools with both solvable and unsolvable tasks to test whether models recognize their limitations. GeoAnalystBench \citep{zhang2025geoanalystbench} evaluates Python code generation for spatial analysis workflows, addressing uncertainty about LLMs' actual capabilities in automating GIS tasks. Systems such as LLM-Geo \citep{li2023llmgeo} and GeoGPT \citep{zhang2024geogpt} demonstrate end-to-end geospatial problem solving through tool integration.

\paragraph{Mathematical reasoning.}
Mathematical reasoning benchmarks provide context for understanding computational capabilities. GSM8K \citep{cobbe2021gsm8k} tests grade-school math word problems requiring 2--8 step solutions, where frontier models now exceed 90\% accuracy. The MATH dataset \citep{hendrycks2021math} contains 12,500 competition-level problems with step-by-step LaTeX solutions, though recent models like GPT-o1 achieve 94.8\%. More challenging benchmarks include AIME \citep{aime2024} for competition mathematics and FrontierMath \citep{glazer2024frontiermath} for research-level problems where models score below 2\%. GeoEval \citep{zhang2024geoeval} specifically targets geometry, but focuses on planar rather than spherical geometry.

\paragraph{Positioning of \dsName.}
\dsName differs from existing benchmarks in three key aspects. First, unlike geographic knowledge benchmarks that test factual recall (country statistics, place attributes), \dsName requires \textit{coordinate-level reasoning}: computing distances, bearings, and spatial relationships from raw GPS coordinates. Second, unlike tool-augmented benchmarks that evaluate API orchestration, \dsName tests \textit{intrinsic capabilities} encoded in model parameters, revealing training gaps that external tools would mask. Third, unlike mathematical reasoning benchmarks focused on planar geometry, \dsName uniquely tests \textit{geodetic computation}: Haversine distances, great-circle interpolation, and spherical polygon areas---formulas essential for geographic applications but absent from existing benchmarks.
